At the end of last section, we stated that information from ganglion cells extend to the Lateral Geniculate Nucleus (LGN), where information is relayed and organized so that the cortex can interpret it. Organization makes left visual field sent to right hemisphere, right field to left hemisphere (Figure~\ref{fig:vision:optic-chiasm}).

\begin{figure}
  \begin{center}
    \includegraphics[width=0.7\textwidth]{visual-fields-lgn}
    \caption{Visual field mapping in LGN.}
    \label{fig:vision:optic-chiasm}
  \end{center}
\end{figure}

The portion of the cortex that is involved with visual processing has been estimated to about 30\%.

It has been studied and areas have been labelled due to their function.

V1, 
The primary visual cortex (V1) is the best-studied visual area in the brain.
correspondence between a given location in V1 and in the subjective visual field is very precise
 fovea in the retina, a large portion of V1 is mapped to the small, central portion of visual field
neuronal responses can discriminate small changes in visual orientations, spatial frequencies and colors. 
described as edge detection. 

V2
connections from v1 sends to v3,v4,v5 back to b1
more complex shapes than v1

V3
controversy still exists regarding the exact extent of area V3,
 it may contain a complete visual representation
  may play a role in the processing of global motion
  
V4
strong feedforward input from V2 and sending strong connections to the PIT
receives direct inputs from V1, especially for central space
it has weaker connections to V5 and dorsal prelunate gyrus (DP).
show strong attentional modulation
orientation, spatial frequency, and color.
tuned for object features of intermediate complexity, like simple geometric shapes
V4 is not tuned for complex objects such as faces, as areas in the inferotemporal cortex are

V5
Visual area V5, also known as visual area MT (middle temporal), is a region of extrastriate visual cortex that is thought to play a major role in the perception of motion, the integration of local motion signals into global percepts, and the guidance of some eye movements.


V6
sharp selectivity for the orientation of visual contours, and preference for long, uninterrupted lines covering large parts of the visual field
strongly responsive to low spatial frequency components of an image, and respond poorly to the motion of textured pattern
